\begin{lemma}
Let \(M(n)=\noderef{TwoBiteNaturalReducedVertexCount}(n) = \lfloor n/\log^2 n \rfloor\), and let \(\beta = 1/2\). For any \(\eta>0\), under the two-bite law, the probability of the simultaneous red and blue fiber-size concentration event
\[
\big(\forall r\in V_R, \noderef{WithinRelativeError}(\eta, \log^2 n, |\noderef{RedFiber}(\omega,r)|)\big) \land
\big(\forall b\in V_B, \noderef{WithinRelativeError}(\eta, \log^2 n, |\noderef{BlueFiber}(\omega,b)|)\big)
\]
tends to \(1\) as \(n\to\infty\).
\end{lemma}
\begin{proof}
Let \(L=\log n\), \(M=M(n)\), and
\(p=\noderef{TwoBiteEdgeProbability}(1/2,n)\).  By
\(\noderef{TwoBiteNaturalReducedVertexCount}\), together with
\(\noderef{TwoBiteReducedVertexCount}\) and \(\noderef{TwoBiteStretch}\),
\[
M=\left\lfloor \frac n{L^2}\right\rfloor ,
\]
so, after discarding finitely many initial \(n\), we may assume
\[
\frac12\,\frac n{L^2}\le M\le \frac n{L^2},\qquad
\left|\frac nM-L^2\right|\le \frac{\eta}{4}L^2,\qquad
\frac nM\le 2L^2,
\tag{1}
\]
and \(n\le M^2\).  The last inequality is also one of the eventual
conclusions of \(\noderef{TwoBiteNaturalTotalMassOneEventually}\) for
\(\beta=1/2\) and this rounded choice of \(m(n)=M(n)\).  These finite
initial values do not affect \(\noderef{WithHighProbability}\), which is a
limit statement.

We first identify the exact law of one fiber count.  Fix such an \(n\) and a
red base vertex \(r\in \operatorname{Fin}(M)\).  By
\(\noderef{TwoBiteSample}\), an outcome in the two-bite sample space is a
triple
\[
\omega=(G_R,G_B,\pi),
\qquad
\pi:\operatorname{Fin}(n)\hookrightarrow \operatorname{Fin}(M)\times
\operatorname{Fin}(M).
\]
The map \(\pi\) is \(\noderef{TwoBiteEmbedding}(\omega)\).  By
\(\noderef{TwoBiteRedProjection}\), the red projection of a final vertex is
the first coordinate of \(\pi(v)\), and by \(\noderef{RedFiber}\)
\[
|\noderef{RedFiber}(\omega,r)|
=\left|\{v\in \operatorname{Fin}(n):(\pi(v))_1=r\}\right|.
\tag{2}
\]
Thus the random variable in (2) depends only on the injection coordinate
\(\pi\), not on \(G_R\) or \(G_B\).

Let \(A\) be any set of possible values of the count in (2), for example an
upper-tail, lower-tail, or two-sided tail event.  By
\(\noderef{TwoBiteEventProbability}\), the probability of \(A\) is the sum of
\(\noderef{TwoBiteSampleWeight}\) over all triples
\((G_R,G_B,\pi)\) for which the injection \(\pi\) gives a count in \(A\).
The weight factors as
\[
\noderef{TwoBiteSampleWeight}(G_R,G_B,\pi)
=\noderef{GnpGraphWeight}(p,G_R)\,
  \noderef{GnpGraphWeight}(p,G_B)\,
  \noderef{UniformInjectionWeight}(n,M).
\tag{3}
\]
Since membership in \(A\) depends only on \(\pi\), the finite sum factors:
\[
\sum_{\pi:\,X_r(\pi)\in A}\noderef{UniformInjectionWeight}(n,M)
\left(\sum_{G_R}\noderef{GnpGraphWeight}(p,G_R)\right)
\left(\sum_{G_B}\noderef{GnpGraphWeight}(p,G_B)\right).
\]
Both graph sums are \(1\) by \(\noderef{GnpGraphWeightSumOne}\).  Therefore
the marginal probability is exactly
\[
\sum_{\pi:\,X_r(\pi)\in A}\noderef{UniformInjectionWeight}(n,M).
\tag{4}
\]
Because \(n\le M^2\), injections from \(\operatorname{Fin}(n)\) to
\(\operatorname{Fin}(M)\times\operatorname{Fin}(M)\) exist.  Hence
\(\noderef{UniformInjectionWeight}\) is the reciprocal of the number of such
injections, and (4) is precisely the uniform law on ordered samples of \(n\)
distinct cells from the population
\[
U=\operatorname{Fin}(M)\times\operatorname{Fin}(M),\qquad |U|=M^2.
\]
The marked set for the fixed red vertex is the row
\[
R_r=\{r\}\times\operatorname{Fin}(M),
\qquad |R_r|=M.
\]
Equation (2) counts how many of the \(n\) sampled cells lie in this marked
row.  Thus \(X_r=|\noderef{RedFiber}(\omega,r)|\) is a hypergeometric count
with population size \(M^2\), marked-set size \(M\), and sample size \(n\).
Its mean is
\[
\mu=n\frac{M}{M^2}=\frac nM.
\tag{5}
\]

The same reduction holds for a fixed blue vertex
\(b\in \operatorname{Fin}(M)\).  By \(\noderef{TwoBiteBlueProjection}\) and
\(\noderef{BlueFiber}\),
\[
|\noderef{BlueFiber}(\omega,b)|
=\left|\{v\in \operatorname{Fin}(n):(\pi(v))_2=b\}\right|,
\]
and the marked set is the column
\[
B_b=\operatorname{Fin}(M)\times\{b\},
\qquad |B_b|=M.
\]
The graph weights marginalize away exactly as above, so this blue fiber count
has the same hypergeometric law and the same mean \(\mu=n/M\).

We now record the tail estimate used for these hypergeometric variables.
For a sample of size \(s\) without replacement from a population of size
\(N\) with \(K\) marked elements, put \(q=K/N\) and \(\mu=sq\).  The node
\(\noderef{HypergeometricMgfComparison}\) gives, for every real
\(\lambda\),
\[
\mathbb E e^{\lambda Y}
\le (1-q+qe^\lambda)^s
\le \exp(\mu(e^\lambda-1)).
\tag{6}
\]
Let \(0<\delta\le 1/2\).  Applying Markov's inequality to (6) with
\(\lambda=\log(1+\delta)\) gives
\[
\Pr(Y\ge(1+\delta)\mu)
\le \exp\!\left(-\mu((1+\delta)\log(1+\delta)-\delta)\right).
\]
Applying (6) to \(-\lambda\), with
\(\lambda=-\log(1-\delta)\), gives
\[
\Pr(Y\le(1-\delta)\mu)
\le \exp\!\left(-\mu((1-\delta)\log(1-\delta)+\delta)\right).
\]
The two displayed exponents are positive for \(0<\delta\le1/2\), so with
\[
c_\delta=\min\{(1+\delta)\log(1+\delta)-\delta,\,
                (1-\delta)\log(1-\delta)+\delta\}>0
\]
we have
\[
\Pr(|Y-\mu|>\delta\mu)\le 2\exp(-c_\delta\mu).
\tag{7}
\]

Set
\[
\delta=\min\{\eta/4,\,1/4\}.
\]
Then \(0<\delta\le1/2\).  For the red fiber count \(X_r\), (5) and (7) give
\[
\Pr(|X_r-\mu|>\delta\mu)\le 2e^{-c_\delta\mu}.
\tag{8}
\]
If \(|X_r-\mu|\le\delta\mu\), then by (1) and
\(\delta\le\eta/4\),
\[
|X_r-L^2|
\le |X_r-\mu|+|\mu-L^2|
\le \delta\mu+\frac{\eta}{4}L^2
\le \frac{\eta}{4}\cdot 2L^2+\frac{\eta}{4}L^2
\le \eta L^2.
\]
This is exactly
\(\noderef{WithinRelativeError}(\eta,L^2,X_r)\).  Therefore failure of the
red relative-error statement implies the tail event in (8).  Since
\(\mu=n/M\) and (1) gives \(\mu\ge L^2\) for all sufficiently large \(n\),
\[
\Pr\bigl(\neg\noderef{WithinRelativeError}
(\eta,L^2,|\noderef{RedFiber}(\omega,r)|)\bigr)
\le 2e^{-c_\delta L^2}.
\tag{8a}
\]
The identical estimate holds for every blue vertex \(b\), with
\(\noderef{BlueFiber}\) in place of \(\noderef{RedFiber}\).

There are \(M\) red vertices and \(M\) blue vertices.  By the union bound, the
probability that any red or blue fiber violates the required relative-error
bound is at most
\[
4M e^{-c_\delta L^2}.
\tag{9}
\]
Using \(M\le n/L^2\) from (1), the logarithm of the right side is at most
\[
\log 4+\log n-2\log L-c_\delta L^2,
\]
which tends to \(-\infty\) because \(L=\log n\) and \(L^2\) dominates \(L\).
Thus the failure probability tends to \(0\).  Equivalently, the
\(\noderef{TwoBiteEventProbability}\) of the simultaneous red and blue
fiber-size concentration event tends to \(1\), which is the asserted
\(\noderef{WithHighProbability}\) conclusion.
\end{proof}
